\begin{eabstract}
Hedge fund trading is a highly systematic, automated, and game-oriented decision-making activity in financial markets. Its core is not a single price prediction task, but a continuous decision process that involves asset selection, factor analysis, timing identification, trading execution, and final trading decisions. Under long-short competition, leverage constraints, high-dimensional non-stationary time series, and strict execution cost requirements, traditional financial prediction models and single adversarial learning paradigms are prone to insufficient strategy exploration, mode collapse, noise overfitting, and amplified execution deviations. Developing robust, generalizable, and interpretable adversarial learning methods along the real trading process is therefore of theoretical significance and practical value for improving the risk-return performance of hedge funds in complex market environments.

In real trading scenarios, hedge fund decisions are not completed by issuing a single buy-or-sell signal. Instead, they proceed through a layered process involving asset allocation evidence, factor-based explanation, temporal-structure evidence, execution-path constraints, and final strategy integration. Specifically, the asset allocation stage identifies relative allocation value and determines capital weights across multiple assets; the factor modeling stage extracts stable and interpretable information from high-dimensional financial variables; the temporal structure identification stage determines trend patterns and trading timing under non-stationary market conditions; the execution stage controls slippage, market impact, and execution deviations when orders enter the market; and the strategy integration stage organizes directional judgments, structural signals, and execution constraints into buy, sell, or hold actions. Existing methods still face several limitations along this decision chain, including unstable characterization of multi-asset allocation relationships, noise-sensitive factor selection, structural identification affected by local pseudo-patterns, insufficient endogenous control of execution-layer risks, and the lack of unified constraints for final trading decisions. These coupled problems make the practical application of adversarial learning to hedge fund trading a systematic challenge.

To address these problems, this dissertation does not regard hedge fund trading algorithms as a single model optimization problem. Instead, it treats them as a complete trading decision process consisting of portfolio management, high-dimensional financial factor modeling, temporal structure identification, trading execution optimization, and strategy integration. Around this process, this dissertation proposes several multi-agent adversarial learning methods for key links of hedge fund trading, so that each chapter serves a specific problem in the trading decision chain. The main contributions are as follows.

First, for the asset allocation and portfolio management stage, a Group Cross Adversarial (GCA) learning method for portfolio management is proposed. To address the limited strategy exploration space under single-opponent conditions, unstable characterization of multi-asset relative relationships, and portfolio weight generation disturbed by local signals, the method divides agents into several groups and introduces cross-group adversarial interactions, forming an orderly group cross-adversarial game structure. It is used to characterize capital allocation relationships among multiple assets and suppress overfitting along a single game path. Experiments are conducted on multi-instrument futures and equity portfolio trading tasks. The results show that the proposed method significantly outperforms traditional baselines and demonstrates strong robustness in terms of Sharpe ratio and maximum drawdown control, verifying its effectiveness in multi-asset capital allocation and portfolio risk control.

Second, for the stage of trading-rationale explanation and high-dimensional financial factor modeling, a racetrack competition method for high-dimensional financial factor modeling is proposed. To address the problems that effective information in high-dimensional factor spaces can be masked by noise, model representations can be suppressed by asymmetric competition, and individual models can converge prematurely, the method introduces dynamic competition among multiple heterogeneous models under a unified scoring rule. It further constructs a bidirectional knowledge transfer logic in which mature models provide stable representations to emerging models, while emerging models provide differentiated feedback constraints to mature models. In this way, the method selects and strengthens financial factors with stable explanatory power and delays premature convergence of any single model. Validation experiments are conducted in complex coupled-factor simulation environments and on real high-frequency data. The results show that the racetrack competition method significantly improves the robustness of high-dimensional factor representations and effectively alleviates information loss during factor extraction.

Third, for the stage of trading-timing judgment and temporal structure identification, a reverse student-to-teacher distillation learning method for temporal structure identification is proposed. To address the difficulty of quantifying subjective wave patterns, the tendency of the conventional ``strong-teaching-weak'' paradigm to fall into local optima under non-stationary noise, and the difficulty of determining the trading meaning of structural signals, the method combines quaternary fractal events to quantitatively reconstruct Elliott Wave Theory and designs a reverse distillation path that transfers failure experience from weak models to strong models. This encourages the best-performing model to align with the intra-group distribution and suppress noise overfitting. Experiments are conducted under market conditions with complex trend structures. The results show that the proposed method outperforms mainstream baseline models in structural identification performance, verifying its effectiveness in filtering high-frequency white noise, preserving trend structure information, and supporting trading timing identification.

Fourth, for the stage of execution-path optimization and trading execution control, a twin adversarial self-regeneration method for trading execution optimization is proposed. To address volume prediction deviations, execution path imbalance, slippage risk, and adversarial collapse after trading signals enter the real execution stage, the method uses a ``copy-pair-differentiate'' self-repair logic to automatically trigger cross-group repair procedures when intra-group adversarial collapse occurs, thereby internalizing market impact risk into optimizable mathematical objectives. Using a volume-weighted average price execution strategy as the validation object, experiments are conducted on real trading data. The results show that the method has clear advantages in execution slippage control and provides a feasible collapse recovery solution for fine-grained trading execution optimization.

Fifth, for the stage of final trading decision and strategy integration, a Partial-order Collaborative Adversarial (PCA) learning method for strategy integration is proposed. To address hierarchical game relationships, mode collapse, and the difficulty of unifying multiple constraints in final trading decisions, the method adopts an asymmetric configuration to simulate interactions among market participants with different capital sizes. It further interprets upward, downward, and neutral classification outputs as buy, sell, or hold actions, and connects asset allocation, factor evidence, structural signals, and execution constraints at the levels of problem definition, model interpretation, and evaluation criteria. Comparative validation is conducted across different asset classes with asymmetric game characteristics. The results show that the proposed method performs strongly in cross-asset generalization and significantly improves the robustness of strategy integration.

In summary, this dissertation proposes several multi-agent adversarial learning methods around the hedge fund trading decision chain. Chapter 2 uses GCA to support capital allocation across multiple assets. Chapter 3 develops a dynamic mechanism for identifying effective financial factors. Chapter 4 identifies temporal structures for trend and timing judgment. Chapter 5 optimizes trading execution during order implementation. Chapter 6 uses PCA to organize directional classification into final trading decisions and strategy integration. These studies proceed from asset allocation, factor selection, structural identification, and execution optimization to final strategy integration, and they can effectively address modeling challenges posed by high-dimensional non-stationary time series, complex game structures, and extreme market conditions, providing more robust, efficient, and interpretable decision support for hedge fund trading.
\end{eabstract}

\ekeywords{Hedge Fund Trading, Multi-Agent Adversarial Learning, Group Cross Adversarial, Partial-order Collaborative Adversarial, Strategy Integration}
